\section{Gap individuato e posizionamento della tesi}
\label{sec:gaps}

I capitoli precedenti hanno descritto un panorama di ricerca maturo. La teoria dei giochi cooperativi offre il riferimento canonico per la ripartizione dei benefici \cite{shapley1953, moncecchi2020}; un insieme di approssimazioni scalabili ne ha reso fattibile il calcolo anche per comunità grandi \cite{castro2009, cremers2023, alonsopedrero2024, ferrucci2025}; gli approcci semplificati restano attraenti per la loro trasparenza verso utenti non esperti \cite{kulmala2021, fioriti2023}; e un gruppo di indicatori di equità, importati da altre discipline, permette di valutare le ripartizioni risultanti (Sezione~\ref{sec:2.3}). Il problema aperto non è più \emph{come calcolare} una ripartizione equa, ma \emph{come scegliere} fra ripartizioni che possono tutte essere calcolate.

\subsection*{Il bisogno di ricerca da cui parte il lavoro}

L'analisi assiomatica più recente valuta otto metodi di ripartizione rispetto a sette proprietà desiderabili e conclude che \textit{``no computationally efficient CAM satisfies all desirable properties simultaneously, highlighting structural tensions between fairness axioms, incentive compatibility, and ease of implementation''} \cite{gonzalezAsenjo2025}. Per una famiglia di metodi la tensione è un teorema: \textit{``this trade-off is inherent to the entire price-based family''} \cite{gonzalezAsenjo2025}. Gli autori dichiarano anche dove il risultato si ferma. Il loro modello è \textit{``intentionally minimal''}, con \textit{``known and deterministic consumption and generation within each time slot''}, e fra gli sviluppi futuri indicano l'estensione del quadro a contesti con \textit{``household flexibility, or multi-period coupling''}, che \textit{``would help evaluate the robustness of CAM properties under more realistic conditions''} \cite{gonzalezAsenjo2025}. Dei quattro sviluppi elencati, questa tesi riprende il terzo e, per l'uso di profili realistici, il quarto; i primi due, sulla caratterizzazione formale del nucleolo e sugli algoritmi efficienti per calcolarlo, non sono perseguiti. Un confronto data-driven degli stessi autori va nella stessa direzione: \textit{``no single cost allocation method dominates across all dimensions and configurations''}, e \textit{``formal property compliance alone is insufficient to assess real-world performance and must be complemented by scenario-specific analysis''} \cite{gonzalezAsenjo2026}.

Da queste affermazioni restano aperte due domande. Le incompatibilità sono dimostrate per i metodi basati su prezzi e per un singolo intervallo di tempo; se tensioni comparabili compaiano per i meccanismi semplificati, performance-based e game-theoretic, quando il tempo è esplicito e i membri differiscono per il sincronismo delle loro immissioni e dei loro prelievi, per quanto ci risulta non è stato quantificato. Inoltre gli assiomi sono soddisfatti o violati, mentre la scelta di una regola richiede di sapere \emph{di quanto} essa rinuncia a una proprietà per ottenerne un'altra.

Rispondere alla seconda domanda richiede un metro comune, e la letteratura indica più volte che manca. Dynge e Cali osservano che gli indicatori di equità in uso sono stati \textit{``adapted for use in the context of local electricity markets''} da altre discipline e che, prima del loro studio, \textit{``no research has been found addressing or comparing the performance of these indicators''}; concludono che \textit{``additional investigation is needed to refine fairness indicators, ensuring their effectiveness across diverse market models''} \cite{dynge2025}. Gasca et al.\ avvertono che \textit{``fairness metrics derived from other fields are being applied indiscriminately to ECs' context''} \cite{gasca2024}, Volpato et al.\ indicano fra gli sviluppi futuri la necessità di \textit{``identify parameters that can quantitatively compare the different benefit distributions''} \cite{volpato2024}, e per Cremers et al.\ \textit{``there are many other fairness concepts that could be explored in energy applications, and it would be relevant to compare their outcome to Shapley values''} \cite{cremers2023}. Una rassegna recente segnala la mancanza di \textit{``conceptual alignment between performance metrics''}, usate insieme \textit{``without a systematic articulation of how these criteria interact or conflict''}, e conclude che efficienza, equità, scalabilità e inserimento regolatorio \textit{``are not orthogonal dimensions. Improving one dimension often weakens another''} \cite{chatel2026}. Lo stesso vocabolario non è condiviso: l'equità è \textit{``the equal allocation of costs/benefits for all users, regardless of their contribution''} per Gasca et al.\ \cite{gasca2024}, mentre per Dynge e Cali \textit{``Equity commonly refers to rewarding individuals proportionally to individual contribution''} \cite{dynge2025}.

Che la scelta della regola conti anche quando il valore da ripartire è fissato ex post è sostenuto, in un modello stilizzato con due agenti e senza profili orari, da Awerkin et al., per i quali la regola di ripartizione \textit{``is not a mere redistribution device but a determinant of the equilibrium investment pattern''} e la ripartizione del beneficio collettivo è \textit{``a regulatory issue in its own right''} \cite{awerkin2026}. Questa tesi non studia come i membri rispondono alla regola; misura che cosa ciascuna regola premia.

I confronti esistenti sono circoscritti. Il più vicino è quello di Couraud et al., che valutano quattro meccanismi di mercato locale (pro-rata, due forme di glass-filling e una doppia asta) su quattro nozioni di equità in 250 comunità simulate di venti membri, e concludono che \textit{``there is not one distribution approach that addresses all types of fairness''} \cite{couraud2025}. Il loro risultato è letto da grafici, riguarda meccanismi che operano a un prezzo locale uniforme, varia solo la penetrazione del fotovoltaico e non include né regole cooperative, né stabilità di coalizione, né forza incentivante; l'equità basata sul bisogno è esclusa per mancanza di dati. Con incentivi regolati sull'energia condivisa, gli studi che confrontano più di un meccanismo poggiano su un singolo caso: tre membri in \cite{volpato2023}, cinque in \cite{trevisan2024}, dieci in \cite{fioriti2023} e un solo comune in \cite{campagna2024}; lo studio che varia taglia, quota di prosumer e diversità degli utenti su 64 casi applica un solo meccanismo \cite{ferrucci2025}. Gasca et al.\ incrociano dieci approcci e quattro indici e concludono che \textit{``the findings are sensitive to user assets and profile changes''}, ma anche che \textit{``the same conclusions apply to other types of ECs with diverse user types''} \cite{gasca2024}: una generalità che il loro disegno, una sola comunità di sette famiglie osservata per un mese, non verifica.

A nostro avviso il gap è quindi il seguente. Il test se il verdetto sui meccanismi dipenda dalla nozione di equità è stato condotto per quattro meccanismi di mercato locale e quattro nozioni \cite{couraud2025}. Ciò che i lavori esaminati non mostrano è come meccanismi di famiglie diverse si ordinino sugli stessi dati con un insieme comune di indicatori che includa, oltre a uguaglianza e proporzionalità, la \emph{stabilità di coalizione} e la \emph{forza incentivante}; quanto si perde su una nozione per guadagnare su un'altra, come grandezza e non come grafico; e come l'ordinamento cambi quando taglia, penetrazione dei prosumer e composizione variano in modo sistematico, quando il valore da ripartire è un trasferimento regolato sull'energia condivisa oraria.

\subsection*{Contributo della tesi e sue estensioni}

La tesi affronta questo bisogno con un modello tecnico-economico di simulazione a risoluzione oraria, in cui i meccanismi sono valutati nelle stesse condizioni operative e con gli stessi indicatori. Da questo nucleo il lavoro si estende in tre direzioni, ciascuna motivata dalla letteratura.

La prima direzione riguarda l'insieme dei meccanismi. Il confronto data-driven copre tre delle sei famiglie della Sezione~\ref{sec:2.2} e lascia il quadro aperto a \textit{``methods not considered in this paper''} \cite{gonzalezAsenjo2026}, e Gasca et al.\ elencano fra gli sviluppi futuri \textit{``additional types of comparisons''} \cite{gasca2024}. Poiché le tensioni dimostrate riguardano i metodi basati su prezzi, la domanda aperta riguarda proprio i metodi che non li usano, fra cui quelli performance-based, progettati per premiare l'allineamento temporale fra immissioni e prelievi \cite{gianaroli2024, giudice2026}. La famiglia basata su prezzi interni e tariffe dinamiche (Sezione~\ref{sec:2.2.4}) fa eccezione: presuppone un prezzo interno, che il quadro italiano non prevede, e non è quindi rappresentata.

La seconda direzione riguarda lo strato di valutazione. Il rispetto formale delle proprietà non basta \cite{gonzalezAsenjo2026}, e lo strumento per l'analisi complementare non è ben consolidato. Il lavoro impiega quindi un insieme di indicatori e non uno solo, tratta il loro accordo come un risultato da osservare, e riporta a parte gli indicatori definiti come distanza da un'allocazione di riferimento, il cui valore è per costruzione ottimo per il meccanismo che fornisce il riferimento. Le grandezze calcolate sono analoghe, non identiche, alle proprietà assiomatiche: il lavoro le quantifica e non verifica il teorema.

La terza direzione riguarda le configurazioni. L'esigenza di variare la composizione ricorre fra gli sviluppi futuri: comunità con \textit{``a higher number of members (e.g., tens or more)''} \cite{volpato2024}, \textit{``the different EC compositions (number of members in the UT groups)''} \cite{minutoLanzini2022}, \textit{``the critical set-up''} di una comunità oltre il quale i meccanismi di ridistribuzione restano applicabili \cite{fina2024}, e \textit{``the number and composition of community members''} \cite{giudice2026}. Il confronto è quindi ripetuto su configurazioni che variano il numero di membri, la penetrazione dei prosumer e la composizione per categoria di membri.
% TODO: la composizione per categoria è prevista ma non ancora eseguita (tempi di calcolo, README Sezione 12): aggiungere le configurazioni prima della consegna o cancellare l'ultima clausola.

\subsection*{Domanda di ricerca e contesto empirico}

Il gap e le sue estensioni convergono in un'unica domanda di ricerca:

\begin{quote}
\emph{Quando i profili orari di un anno sono ripartiti con meccanismi di famiglie diverse, (i) l'ordinamento dei meccanismi cambia con la nozione di equità adottata (uguaglianza, proporzionalità al contributo fisico, stabilità di coalizione, inclusione) e con la taglia, la penetrazione dei prosumer e la composizione per categoria della comunità; (ii) quanto sono grandi i trade-off fra queste nozioni, inclusa la forza incentivante della regola, lungo la frontiera di Pareto; e (iii) quale sottoinsieme minimo di indicatori basta a discriminare i meccanismi?}
\end{quote}

Ogni termine corrisponde a una grandezza calcolata dal modello. L'equità è misurata da indicatori di uniformità, di proporzionalità al contributo fisico, di inclusione e di stabilità di coalizione, tutti valutati sulle stesse allocazioni; la stabilità, disponibile fino a diciotto membri, è l'eccesso massimo delle coalizioni bloccanti rispetto al valore da ripartire. La forza incentivante di una regola è misurata ex post, come covarianza trasversale fra la quota ricevuta da ciascun membro e l'allineamento del suo profilo con la comunità, normalizzata su un meccanismo di riferimento; indica quanto una regola premia il sincronismo, non come i membri risponderebbero, una risposta comportamentale che è fuori dal perimetro del modello. L'accordo fra gli indicatori è la correlazione di rango (\(\tau_b\) di Kendall) fra i ranghi che inducono, la dipendenza dalla configurazione è la frequenza con cui l'ordine di una coppia di meccanismi si inverte lungo gli assi di taglia e penetrazione, e il sottoinsieme minimo si ottiene raggruppando gli indicatori in base alla loro distanza di rango e tenendo un rappresentante per gruppo. La sostenibilità economica della partecipazione, calcolata come valore attuale netto dell'investimento di ciascun membro dove le quote di investimento sono specificate, è trattata come analisi complementare.

L'analisi è svolta nel quadro regolatorio italiano, che è il contesto empirico del lavoro e non la sua motivazione. Esso è un'istanza operativa e documentata di un regime in cui il valore da ripartire è un trasferimento regolato e non un arbitraggio: la quantità condivisa è il minimo orario fra immissioni e prelievi aggregati e la remunerazione unitaria è comune a tutti i membri, così che il canale della dispersione dei prezzi, che il confronto data-driven individua come motore dei suoi risultati, è assente, e i membri differiscono solo per l'allineamento temporale dei profili e, dove la norma lo prevede, per categoria giuridica \cite{MASE:Decreto:CER:2023, GSE:CACER:RegoleOperative:2025}. È il caso più semplice in cui porre la domanda. L'ammissibilità di ogni allocazione è verificata ex post rispetto alla soglia del quadro normativo e riportata.
% TODO: verificare questo paragrafo sul testo del decreto e sulle Regole Operative GSE (soglia sul premio, fattore di riduzione, esenzione delle persone fisiche); sulle configurazioni esistenti la soglia non morde (README Sezione 7.4): o la si controlla su tutte le 28 configurazioni, o si riduce l'ultima frase a una nota.

La tesi risponde alla domanda con un modello di simulazione organizzato in due strati. Lo strato di valorizzazione non ha gradi di libertà: a partire dai profili orari dei membri determina il valore da ripartire e i vincoli che agiscono su di esso. Lo strato di allocazione è l'oggetto dello studio e implementa sedici meccanismi di cinque delle sei famiglie della Sezione~\ref{sec:2.2}, tutti ripartiscono lo stesso valore totale e sono valutati con gli stessi indicatori. Il risultato atteso non è l'individuazione di un meccanismo ottimo, ma la struttura dell'interazione fra equità, stabilità e forza incentivante: le configurazioni in cui l'ordinamento si inverte e quelle in cui gli indicatori concordano o divergono. Poiché diversi indicatori sono massimizzati per costruzione da un meccanismo specifico, la grandezza di interesse è la perdita che gli altri meccanismi subiscono sugli indicatori restanti, non l'identità del vincitore su ciascuno.
% TODO: controllare il numero di meccanismi sull'elenco finale (il sedicesimo gira su dati segnaposto e non deve produrre risultati).
